\section{Background}
\label{sec:background}

\subsection{The Regression Tsetlin Machine}

A RegressionTM maintains a set of clauses, each a conjunction of literals over the
Booleanized input features, together with integer weights. A prediction is formed by
summing clause votes and mapping the sum, through a resolution constant $T$, onto the
target range. Each clause is governed by a team of Tsetlin Automata whose states
determine which literals are included; learning proceeds by nudging these automata via
two feedback types.

\subsection{Type~I and Type~II feedback}
\label{sec:feedback}

TM learning is driven entirely by two clause-feedback operations, dispatched per training
sample according to the sign of the prediction error
$\texttt{prediction\_error} = \texttt{pred\_y} - \texttt{encoded\_Y}$:

\begin{itemize}
  \item \textbf{Type~I feedback} (fired when the prediction is \emph{too low},
    $\texttt{prediction\_error} < 0$) reinforces clauses and, in the regression setting,
    pairs with a weight increment---it \emph{raises} the prediction. Type~I includes
    stochastic literal inclusion and is the mechanism by which the model grows its
    representation.
  \item \textbf{Type~II feedback} (fired when the prediction is \emph{too high},
    $\texttt{prediction\_error} > 0$) erodes clauses toward inactivity and pairs with a
    weight decrement---it \emph{lowers} the prediction. Type~II is \emph{self-limiting}:
    it can only drive clause contributions down (bounded at zero), so it cannot run away.
\end{itemize}

The magnitude of each update is set stochastically through
$\texttt{update\_p} = (\texttt{magnitude}/T)^2$, where $\texttt{magnitude}$ is normally
the absolute prediction error. This asymmetry between the two feedback types---one
unbounded and expansive, one bounded and contractive---is central to the stability
results below.
